\documentclass[12pt]{beamer}

\usepackage[language=british,
            sortlocale=en-GB,
            backend=biber,
            natbib=true,
            dashed=false,
            citestyle=authoryear-comp,
            bibstyle=authoryear,
            doi=false,
            useprefix=true,
            sorting=ynt,
            maxbibnames=9,
            maxcitenames=2]{biblatex}

\addbibresource{legenda.bib}
 \renewcommand\nameyeardelim{\space}
 \renewcommand\postnotedelim{\addcolon\space}
 \DeclareFieldFormat{pages}{#1}
 \DeclareFieldFormat{postnote}{#1} % no postnote prefix in "normal" citation commands
 \AtEveryBibitem{%
     \clearfield{labelmonth}%
     \clearfield{issn}%
 }

%%% no navigation icons

\setbeamertemplate{navigation symbols}{}

%% stepwise uncovering

\beamerdefaultoverlayspecification{<+->}

%%% other colors

\definecolor{camcosdark}{RGB}{0,145,193}
\definecolor{camcosblue}{RGB}{0,163,218}
\setbeamercolor*{section in head/foot}{parent=palette quaternary,fg=white,bg=camcosblue}%!90!black}
\setbeamercolor*{subsection in head/foot}{parent=palette secondary,fg=white,bg=camcosblue}
\setbeamercolor{structure}{bg=white,fg=black}
\setbeamercolor*{palette primary}{use=structure,bg=white,fg=black}
\setbeamercolor*{palette secondary}{use=structure,fg=white,bg=camcosblue}%!90!black}
\setbeamercolor*{palette tertiary}{use=structure,fg=white,bg=camcosblue}%!90!black}
\setbeamercolor*{palette quaternary}{use=structure,fg=white,bg=camcosblue}%!90!black}

\setbeamercolor*{alerted text}{fg=camcosblue}%!90!black}
\setbeamercolor*{part in toc}{fg=camcosblue}
\setbeamercolor*{section in toc}{fg=camcosblue}
\setbeamercolor*{item}{fg=camcosblue}

\setbeamertemplate{headline}
{%
\leavevmode
\hbox{\begin{beamercolorbox}[wd=1.00\paperwidth, ht=3.0ex, dp=1.5ex, right, rightskip=.5cm]{section in head/foot}%
        \usebeamerfont{title in head/foot}\scriptsize{\insertsection{}}
    \end{beamercolorbox}

    \begin{beamercolorbox}[wd=.00\paperwidth, ht=2.5ex, dp=1ex, left, leftskip=.5cm]{subsection in head/foot}%
        \usebeamerfont{title in head/foot}%\insertsubsection
    \end{beamercolorbox}
    }
}

\setbeamertemplate{footline}
{%
  \leavevmode%

  \begin{beamercolorbox}[wd=1\paperwidth,ht=2.5ex,dp=1.5ex,leftskip=.3cm plus1fill,rightskip=.3cm]{}%
    \usebeamerfont{title in head/foot}\insertframenumber/\inserttotalframenumber%
  \end{beamercolorbox}%
}

\setbeamertemplate{partpage}
{%
  \begin{centering}
    {\usebeamerfont{part name}\usebeamercolor[bg]{part name}\partname~\insertpartnumber}
    \begin{beamercolorbox}[sep=16pt,center]{part title}
      \usebeamerfont{part title}\insertpart\par
    \end{beamercolorbox}
  \end{centering}
}

\defbeamertemplate*{title page}{customised}[1][]
{%
    \usebeamerfont{title}\alert{\inserttitle}\par
    \usebeamerfont{subtitle}{\large\insertsubtitle}\par
    \bigskip
    \hrule
    \usebeamerfont{author}\insertauthor\par
    \usebeamerfont{institute}\small\insertinstitute\par
    \bigskip
    \hrule
    \usebeamerfont{date}{\normalsize\insertdate}\par
    \usebeamercolor[fg]{titlegraphic}\inserttitlegraphic{}
}

%%%

\setbeamertemplate{caption}[numbered]
\setbeamertemplate{caption label separator}{:}
\setbeamercolor{caption name}{fg=normal text.fg}

\usepackage{amssymb,amsmath}
\usepackage{mathtools}
\usepackage{fixltx2e} % provides \textsubscript

\usepackage{mathspec,xunicode,fontspec,xltxtra}
\defaultfontfeatures{Mapping=tex-text,Scale=MatchLowercase}
\setmainfont[BoldFont = Source Sans Pro Semibold, BoldItalicFont = Source Sans Pro Semibold Italic]{Source Sans Pro}
\setsansfont[BoldFont = Source Sans Pro Semibold, BoldItalicFont = Source Sans Pro Semibold Italic]{Source Sans Pro}
\setmonofont[Scale=0.8]{Source Code Pro}
\setmathfont(Digits,Latin,Greek)[Lowercase=Regular,Uppercase=Regular]{Clear Sans}

\usepackage{microtype}

% linguistics
\usepackage{tikz-qtree}
\usetikzlibrary{positioning,decorations.pathmorphing,decorations,trees,fit}
\tikzstyle{arrow} = [thin, <-, below left, bend right, font = \scriptsize]
\tikzset{>=latex}
\tikzset{every tree node/.style={align=center,anchor=north}}
\tikzset{execute at begin node=\vphantom{A}}
\newcommand*\circled[1]{\tikz[baseline=(char.base)]{
    \node[shape=circle, draw, inner sep=2pt] (char) {#1};}}

\usepackage{pgfplots}

\usepackage{expex}
\lingset{aboveglftskip=-1ex,
        belowglpreambleskip=-1ex}

\resetcountonoverlays{excnt}

\usepackage{url}

\usepackage{graphicx}
\makeatletter
\def\maxwidth{\ifdim\Gin@nat@width>\linewidth\linewidth\else\Gin@nat@width\fi}
\def\maxheight{\ifdim\Gin@nat@height>\textheight0.8\textheight\else\Gin@nat@height\fi}
\makeatother
% Scale images if necessary, so that they will not overflow the page
% margins by default, and it is still possible to overwrite the defaults
% using explicit options in \includegraphics[width, height, ...]{}
\setkeys{Gin}{width=\maxwidth,height=\maxheight,keepaspectratio}

% Comment these out if you don't want a slide with just the
% part/section/subsection/subsubsection title:
\AtBeginPart{%
  \let\insertpartnumber\relax
  \let\partname\relax
  \frame{\partpage}
}
\AtBeginSection{%
  \let\insertsectionnumber\relax
  \let\sectionname\relax
  \frame{\sectionpage}
}
\AtBeginSubsection{%
  \let\insertsubsectionnumber\relax
  \let\subsectionname\relax
  \frame{\subsectionpage}
}

\usepackage[normalem]{ulem}
% avoid problems with \sout in headers with hyperref:
\pdfstringdefDisableCommands{\renewcommand{\sout}{}}

\setlength{\parindent}{0pt}
\setlength{\parskip}{6pt plus 2pt minus 1pt}
\setlength{\emergencystretch}{3em}  % prevent overfull lines

\setcounter{secnumdepth}{1}

\title{Agreement with possessed direct objects in Hungarian}

\subtitle{An experimental approach}

\author{\textsuperscript{a}András Bárány, \textsuperscript{b}Ádám Szalontai}

\date{SinFonIJA 8, Ljubljana, 25 September 2015}

\institute{\textsuperscript{a}ReCoS/DTAL, University of Cambridge,\\
        \textsuperscript{b}Research Institute for Linguistics,
        Hungarian Academy of Sciences}

\titlegraphic{%
    \begin{tikzpicture}[overlay]
        \node at (8.2, -2.5) {\includegraphics[height=2.5cm]{erc.eps}};
    \end{tikzpicture}
}

\begin{document}

\frame{%
\thispagestyle{empty} % no page number here
  \titlepage{}
}

\begin{frame}{Overview}
\tableofcontents[hideallsubsections]
\end{frame}

\section{Intro: Object agreement and possessed NPs}\label{case-and-case}

\begin{frame}{Object agreement in Hungarian}

\pex
    \a
        \begingl
            \gla    Mari lát egy kutyá-t.//
            \glb    M.      see.\textsc{3sg} a   dog-\textsc{acc}//
            \glft   ‘Mari sees a dog.’//
        \endgl
    \pause
    \a
        \begingl
            \gla    Mari lát-ja a kutyá-t.//
            \glb    M.   see-\textsc{\alert{3sg.obj}} the dog-\textsc{acc}//
            \glft   ‘Mari sees the dog.’//
        \endgl
\xe

    \begin{itemize}

        \item Roughly sensitive to definiteness of the DO

        \item Possessed direct objects are an exception

        \item They can be indefinite yet require agreement

        \item Standard Hungarian: all possessed DOs agree

        \item But: dialectal variation?

    \end{itemize}

\end{frame}

\begin{frame}{This presentation}

    \begin{block}{Claims}

        \begin{itemize}

        \item Variation is difficult to find

        \item Speakers treat all possessed direct objects in the same way

        \item Agreement is not semantic, but sensitive to \alert{syntactic structure}

        \end{itemize}

    \end{block}

    \begin{block}{Methods}

        \begin{itemize}

            \item Series of online surveys

            \item Acceptability judgments with and without contexts

            \item Forced choice tests with and without contexts

        \end{itemize}
    \end{block}

\end{frame}

\section{The data}

\begin{frame}{Object agreement}

    \begin{itemize}

        \item Definiteness often a good predictor; not with possessed DOs

        \item Indefinite possessed DOs require agreement

    \end{itemize}

\pex
    \a
        \begingl
            \gla    Mari {lát /} \ljudge*lát-ja egy kutyá-t.//
            \glb    M.   see.\textsc{3sg} see-\textsc{3sg.obj} a   dog-\textsc{acc}//
            \glft   ‘Mari sees a dog.’//
        \endgl
    \pause
    \a
        \begingl
            \gla    Mari {lát-ja /} \ljudge\%lát egy kutyá-m-at.//
            \glb    M.   see-\textsc{\alert{3sg.obj}} see.\textsc{3sg} the dog-\textsc{\alert{1sg.poss}-acc}//
            \glft   ‘Mari sees a dog of mine.’//
        \endgl
\xe

\end{frame}

\begin{frame}{Structures}

\ex
    \begin{tikzpicture}[baseline]
        \Tree   [.DP
                    \node (specdp) {poss.\textsc{nom}};
                    [.{D$'$}
                        D
                        [.PossP
                            \node (t) {\sout{poss.\textsc{nom}}};
                                [.{Poss$'$}
                                    [.NP
                                        N
                                    ]
                                    [.Poss
                                        -\textsc{poss.$\phi$}
                                    ]
                                ]
                        ]
                    ]
                ]
        \draw [arrow, ->, bend left] (t) to (specdp.south);
    \end{tikzpicture}
\xe

\end{frame}

\section{Surveys}

\begin{frame}{Acknowledgements}

\begin{center}
    \Large \textbf{Thank you!}
\end{center}

The first author's research is supported by ERC Grant No. 269752
\enquote{Rethinking Comparative Syntax}. Thanks to: Theresa Biberauer, Anders
Holmberg, Georg Höhn, Ian Roberts, Michelle Sheehan, and Jenneke van der Wal.

\end{frame}

\begin{frame}[allowframebreaks]{References}

\small
\setlength{\parindent}{-0.2in} \setlength{\leftskip}{0.2in}
\setlength{\parskip}{3pt}
\newrefcontext[sorting=nyt]
\printbibliography
\end{frame}

\end{document}
